% Chapter 1: Executive Summary / Introduction
% BSU Engineering Portal Graduation Project

\chapter{Executive Summary / Introduction}
\label{chap:introduction}

\section{Problem Statement}
\label{sec:problem-statement}

The Faculty of Engineering at Beni-Suef University (BSU) operates across multiple academic departments---including Electrical Engineering (with EPM and ECE specializations), Mechanical, Civil, and Architecture Engineering---serving thousands of students across five academic levels (Preparatory through Fourth Year). Historically, academic operations relied on fragmented manual workflows characterized by:

\begin{itemize}
    \item \textbf{Excel-based grade management}: Staff Affairs and Student Affairs personnel maintained disjointed spreadsheets for student registration, course assignments, and grade tracking. Version control was non-existent; data inconsistency across departments was the norm.
    \item \textbf{Ad-hoc communication}: Lecture schedules, exam dates, and announcements propagated through informal channels (WhatsApp groups, email chains), leading to information asymmetry between students and faculty.
    \item \textbf{No audit trail}: Administrative actions---batch student uploads, grade modifications, doctor reassignments---left no traceable record, complicating accountability and regression analysis.
    \item \textbf{Paper-heavy processes}: Attendance tracking, quiz administration, and certificate generation required physical document handling, introducing latency and human error.
\end{itemize}

The bottlenecks and disjointed nature of this manual workflow are illustrated in Figure~\ref{fig:current_state}.

\begin{figure}[H]
    \centering
    \includegraphics[width=0.85\textwidth]{images/manual_workflow_1.png}
    \caption{Diagram showing the manual workflow: Staff upload Excel sheets $\rightarrow$ email to department heads $\rightarrow$ manual data entry $\rightarrow$ printed reports.}
    \label{fig:current_state}
\end{figure}

\section{Motivation}
\label{sec:motivation}

The inadequacy of the current solution manifests in three dimensions:

\begin{enumerate}
    \item \textbf{Operational inefficiency}: A single semester's grade finalization required 2--3 weeks of cross-departmental reconciliation. Errors discovered post-publication necessitated manual re-computation of cumulative GPAs for entire cohorts.
    
    \item \textbf{Security and integrity risks}: Uncontrolled Excel file distribution meant sensitive student data (national IDs, grades) resided on personal laptops and unsecured cloud storage. No role-based access control existed---any staff member with a spreadsheet could modify any student's record.
    
    \item \textbf{Scalability ceiling}: As enrollment grew (current $
    \approx$ 4,000 active students), the Excel-based approach approached breaking point. Real-time queries (``What is the current attendance rate for Level 2 Power students?'') required hours of manual aggregation.
\end{enumerate}

The COVID-19 pandemic accelerated the urgency for digital transformation. Remote instruction demanded a centralized platform where:
\begin{itemize}
    \item Doctors could upload lecture materials and administer online quizzes
    \item Students could access schedules, submit quiz attempts, and track grades in real-time
    \item Administrators could execute batch operations (student uploads, grade imports) with full audit logging
\end{itemize}

\section{Proposed Solution}
\label{sec:proposed-solution}

\textbf{Release BSU Engineering Portal} is a web-based academic management system designed as a \textit{single source of truth} for all faculty operations. The system adopts a modern three-tier architecture, as depicted in Figure~\ref{fig:architecture}:

\begin{itemize}
    \item \textbf{Backend}: Django 5.2 with Django REST Framework (DRF), providing a secure, stateless API layer configured for cross-origin integration.
    \item \textbf{Database}: MySQL 8.0 for relational data with ACID compliance
    \item \textbf{Frontend}: React 19 with Material-UI v7, delivered via Vite build tooling
    \item \textbf{Infrastructure}: Docker containerization with Docker Compose orchestration
\end{itemize}

\begin{figure}[H]
    \centering
    \includegraphics[width=0.95\textwidth]{images/system_architecture.png}
    \caption{High-level system architecture showing the React frontend, Nginx reverse proxy, Django REST API, and MySQL database, all encapsulated within a Docker network.}
    \label{fig:architecture}
\end{figure}

\subsection{Key Capabilities}

\begin{enumerate}
    \item \textbf{Role-Based Access Control (RBAC)}: Five distinct roles (Admin, Doctor, Student, Student Affairs, Staff Affairs) with fine-grained permissions enforced via custom DRF permission classes.
    
    \item \textbf{Academic Entity Modeling}: Hierarchical data model supporting departments, specializations (e.g., Electrical $\rightarrow$ EPM/ECE), academic years, terms, and levels with strict referential integrity.
    
    \item \textbf{Fault-Tolerant Batch Operations}: High-performance Excel/CSV import pipelines using Pandas. Operations are wrapped in atomic transactions to guarantee database consistency during partial failures, as demonstrated below:
    \begin{lstlisting}[language=Python, caption=Atomic Transaction for Batch Processing (\texttt{staff\_affairs\_views.py})]
for index, row in df.iterrows():
    try:
        with transaction.atomic():
            national_id = str(row['national_id']).strip()
            # Process individual row safely; rollback on error without halting batch
            if not User.objects.filter(national_id=national_id).exists():
                User.objects.create_user(username=national_id, ...)
    except Exception as e:
        errors.append(f"Row {index + 2}: {str(e)}")
    \end{lstlisting}
    
    \item \textbf{Quiz Engine}: Multi-modal assessment support (MCQ, Essay, Image-based questions) with time limits and automatic grading for objective questions.
    
    \item \textbf{Audit Trail}: Immutable logging of all administrative actions via the \texttt{AuditLog} and \texttt{UploadHistory} models, ensuring complete accountability for grade modifications and user management.
\end{enumerate}

\section{Project Scope}
\label{sec:project-scope}

\subsection{In Scope}

The BSU Engineering Portal will deliver:

\begin{itemize}
    \item Complete student lifecycle management (enrollment $\rightarrow$ course assignment $\rightarrow$ grading $\rightarrow$ certificate generation)
    \item Doctor-facing tools: lecture uploads, schedule management, quiz creation, grade entry
    \item Student-facing portal: schedule viewing, material access, quiz participation, grade inquiry
    \item Administrative dashboards: batch uploads, role management, audit log review, system-wide announcements
    \item Containerized deployment suitable for on-premises or cloud hosting
\end{itemize}

\subsection{Out of Scope (Explicitly Excluded)}

The following items are \textbf{not} addressed in this release:

\begin{itemize}
    \item \textbf{Mobile native applications}: The system provides responsive web access only. iOS/Android apps are deferred to future iterations.
    \item \textbf{Payment processing}: Tuition fees, fines, and financial transactions remain handled by the university's existing ERP system.
    \item \textbf{Real-time chat/video}: Synchronous communication (live lectures, chat) is excluded; integration with external platforms (Zoom, Teams) is preferred.
    \item \textbf{Horizontal auto-scaling}: While containerized, the initial deployment targets single-node operations. The portal focuses on administrative workflows, not pedagogical content delivery.
\end{itemize}

\subsection{Success Metrics}

The project will be evaluated against the following criteria, which are also summarized in Table~\ref{tab:success_metrics}:

\begin{itemize}
    \item \textbf{Functional completeness}: All functional requirements (detailed in Chapter 3) pass User Acceptance Testing (UAT).
    \item \textbf{Performance}: Page load times $<$ 2 seconds for 95th percentile; API response latency $<$ 500ms for authenticated queries
    \item \textbf{Security}: Zero critical vulnerabilities (OWASP Top 10) at production release
    \item \textbf{Adoption}: 100\% of faculty staff and $>$ 80\% of enrolled students actively use the system within one academic semester of deployment
\end{itemize}

\begin{table}[H]
    \centering
    \begin{tabular}{|p{0.3\textwidth}|p{0.2\textwidth}|p{0.25\textwidth}|p{0.15\textwidth}|}
        \hline
        \textbf{Metric} & \textbf{Target} & \textbf{Measurement Method} & \textbf{Responsibility} \\
        \hline
        Functional Completeness & 100\% pass rate & Acceptance Testing & QA Team \\
        API Response Time & $<$ 500ms & Load Testing & Backend \\
        UI Load Time & $<$ 2s (95th \%) & Browser DevTools & Frontend \\
        Security Vulnerabilities & Zero Critical & SAST / DAST Tools & DevOps \\
        \hline
    \end{tabular}
    \caption{Success Metrics Summary}
    \label{tab:success_metrics}
\end{table}
